\section{Discussion}
\label{sec:discussion}

\paragraph{The envelope is generic; the rules add the classes.}
The main lesson of the production runs is a negative one.
A long-lived demography alone, without any of the new rules, drives nearly every colony to extinction (98/100 in the long-lived control, 199/200 in the corresponding null model run from the initial condition of \textsf{C1}).
Once natality stops, the colony simply ages out, and the phase-D law Q4 is a consistency check of that demography.
Extinction is therefore \emph{not} a result of the rules, and neither is the rest of the envelope: with crowding
damage and the calibrated attraction $\alpha=4$, the candidate without its four rules has a low peak with free space,
an end of natality near the maximum and no rebound (Table~\ref{tab:res-null}), and, scored at the end of the
plasticity window, late-born cohorts that fail to acquire social competence.
We distinguish what the rules encode, what they add, and what the demography and the initial condition set.

\paragraph{Consistency checks of the design.}
Several patterns follow from the rules by construction; they verify the implementation, but they are not findings.
\begin{itemize}
\item \emph{Transplants (O12, O12b, Q1).} Under R1 an adult cannot recover, and conception and pup survival scale
  with the social state of the parents. Phase-D animals have $s\approx0$ (median 0.000 in every arm with R1;
  Table~\ref{tab:res-trans-s}), so they cannot found a colony, with healthy partners or without them (0/60 and 0/120
  in \textsf{C1}). R1 was motivated by Marsden's transplants \cite{Calhoun1973}, and the cross-transplant arm only
  shows that removing the ratchet restores recovery. The non-trivial part is upstream: that the colony drives
  essentially every adult to $s\approx0$, which O10 and O11 measure and which crowding damage alone largely
  achieves. Adults of intermediate social state
  ($0.2\le s\le0.5$, which exist only around weeks 8--20) found a colony in 5 of 60 trials under R1, against 57 of
  60 without R1 and R2: the ratchet also blocks them, through their fecundity, but where they do found a colony
  their offspring settle at the parents' level ($\bar s\approx0.3$--0.4) and the lineage does not collapse
  (Appendix~\ref{sec:res-transplant}).
\item \emph{Persistence of the isolated class (O13p).} R7 was designed after the first candidate produced only a
  transient isolated class, and its capacity $c_R=2$ equals the isolation threshold of O13.
\item \emph{Senescent decline and no social deaths (O8, A4).} Social mortality is off in every preset.
\end{itemize}
Table~\ref{tab:res-null} gives the pass rate of each essential pattern in the null models; the patterns above are
marked $^\ddagger$ in Table~\ref{tab:observations}, whose caption names the rule that implies each of them.

\paragraph{What the rules add.}
Of the eight essential criteria of Table~\ref{tab:res-null}, five (O4, O5, O6, O7 and O10) pass in at least two of
the four null models and carry no evidence for the rules: they verify that the model has a boom, an end of natality
and an extinction. In particular the low peak with free space is generic: the candidate without its four rules
passes the primary O4 in 200 of 200 runs ($\rho_{\rm pk}=N_{\rm pk}/K_\gamma=0.45$, $f_0=0.74$), and the long-lived
control fails it only through its high peak ($\rho_{\rm pk}=0.85$), not through a lack of free space ($f_0=0.59$).
$K_\gamma=m_c|G|$ is the population at which every cell would sit at the damage threshold, not a physical capacity.
What the rules add, and the null models lack, is two patterns:
\begin{itemize}
\item \emph{Coexisting crowded and isolated deteriorated classes} (O13). They are absent in the long-lived control,
  without crowding damage and without the rules, and present with crowd aversion or refuges alone (1.00 in
  \textsf{C1}$'$, 0.92 in \textsf{C0}); full inheritance without social learning also has them. This is a
  calibrated property, obtained in a nearly sessile regime (below), and it is lost with more movement sub-steps.
\item \emph{The persistence of the isolated class} (O13p), which discriminates against three of the four null models
  but holds by design ($c_R=2$ equals the isolation threshold).
\end{itemize}
The failure of the cohorts born during the boom (O11) is not among them. Scored at sexual maturity (six weeks), as
in the primary evaluator, O11 passes in 1.00 of the \textsf{C1} runs against 0.00 in each null model and needs R1
and R2 (0.71 of the applicable runs and 0.00 without each), but that separation was partly definitional: pups are born with
$s\le w\,s_f\le0.2$, the threshold that O11 uses, and at six weeks they are still inside the plasticity window, so
with pups born above the threshold ($s_{\rm base}=0.25$ or 0.5) O11 failed and O17p dropped to 0.08 and 0.00.
Scored at the end of the window ($a_w=12$ weeks) on the same runs, \textsf{C1} still passes (O11 0.995, O17p 0.990),
the dependence on $s_{\rm base}$ disappears (O17p 0.965 and 0.76), and O11 also passes in 0.81 of the long-lived runs
and in 0.755 of the runs without rules, whose late cohorts fall below 0.2 through accumulated crowding damage. Only
against full inheritance without social learning does it still separate in part (0.21). The failure of the late
cohorts is therefore mostly the crowding damage of the envelope, not a contribution of the rules. Against the
long-lived control and the model without rules, the evidence that the primary outcome O17p carries for the rules
rests on O13, with O13p as a design property; against full inheritance without social learning it rests on O11 at
six weeks and is weaker at the end of the window. Crowd aversion is not needed for either pattern (\textsf{C1}$'$
passes them), and it does not supply the free space at the peak: it reduces it ($f_0=0.16$ with AV, 0.52 without).

\paragraph{The central limitation: demography and initial condition.}
The model has no mortality before senescence, and its default initial condition is a synchronous cohort of 400
healthy four-week-old animals. Both matter for the primary outcome. Without background mortality nobody dies until
old age, so the population sits on a plateau of about 51 weeks and $t_{\rm pk}$ is an arbitrary point of it; the ratio
$(T_{\rm ext}-t_{\rm pk})/t_{\rm pk}$ is 2.1, close to Calhoun's 1.84--2.2, with $t_{\rm pk}=\arg\max N$, but 4.0
from the start of the plateau, so we do not count it as agreement. Calhoun recorded adult deaths during phases B and
C. With a background hazard of 0.003 per week at all ages the plateau shrinks to 17 weeks and O17p falls from 0.995
to 0.55; with 0.008 per week, to 0.06. With initial ages spread over 4--52 weeks O17p is 0.93, but over 4--80 weeks it
is 0.63 (Section~\ref{sec:res-robustness}). In every case the only criterion that fails is O5, and it fails by its
timing clause, not by its level clause. Natality still ends with the population near its maximum (at the 99th
percentile of the births the population is about $0.99\,N_{\rm pk}$), but the deaths bring the start of the plateau
forward ($t'_{\rm pk}$ moves from 34 to 29 weeks at 0.003 per week, medians over 200 runs) while the end of natality
hardly moves (40 to 43 weeks), so $(t_{B99}-t'_{\rm pk})/t'_{\rm pk}$ crosses 0.5 and O5 becomes MARGINAL. This is
the case in 83 of the 90 failing runs at 0.003 per week and in 73 of the 74 failing runs with ages over 4--80 weeks;
the remaining ones, and the 70 runs that fail O5 at 0.008 per week, exceed the timing clause by more than a factor
two. O5 is generic, and the patterns that the rules add (O13 and O13p), as well as O11, pass in every run of these
arms. The
dependence on the demography and on the initial condition is therefore a dependence of the timing of the end of
natality relative to the start of the plateau, a quantity that also rests on the post hoc choice of $t'_{\rm pk}$.
It is not evidence against what the rules add, nor for it, since O5 discriminates against no null model; but it
does mean that O17p $=0.995$ overstates how robustly the full conjunction of patterns is reproduced. A nest period
of three weeks (0.835, borderline) and a random sequential update (1.00) preserve the primary outcome.

\paragraph{The dynamics is a boom and its senescence, a compressed sequence of phases.}
The time course of \textsf{C1} is essentially one transient (Tables~\ref{tab:calib-generations}
and~\ref{tab:calib-chrono}; reference ensemble, 200 runs). The founders deteriorate by crowding during their own adult
life, from $s=0.90$ at the end of the plasticity window to 0.20 at 26 weeks of age. Their offspring, about 70\% of
all births, are born below the deterioration threshold and fail to acquire competence, because competent adults have
become scarce. The net reproductive number falls from 3.8 in the founders to 0.40 in their daughters, so natality
stops within about one generation, and the third and later generations add 3.2\% of the births. The null
models without the rules have a \emph{deeper} boom: with the same attraction and crowding damage, \textsf{C1}
without its four rules has a mean generation of the births of 1.59 against 1.34 and a net reproductive number of 0.66
in the first generation against 0.40, and the long-lived control 1.95 and 1.2 (30 runs each;
Table~\ref{tab:calib-gamma}). The rules, R2 in particular ($-$R2$^*$ gives 1.62), shorten the generational depth
rather than create a long cascade.
The shallowness of the boom is set by how pups are born and socialised: they are born below the deterioration
threshold ($w=0.2$, $s_{\rm base}=0$), conception and pup survival scale with the parents' social state, and social
learning (R2) draws on adults who are already deteriorating, so the net reproductive number drops in the first
generation. A ratio of time scales gives a heuristic comparison with Universe~25, not the cause. Measured on animals
past the plasticity window, so that newborns do not dilute it, $\Gamma=t_{\rm det}/T_{\rm gen}$ is
$\approx1.1$ (1.2 on the founders; about 0.7 if $T_{\rm gen}$ is taken as the mean age of the mothers at birth,
14.3 weeks, so it is of order one under either definition); in Universe~25 the ``death of societal organization'' took the whole of phase~C, about 35
weeks, with a generation time of about 10 weeks, so $\Gamma\approx3.5$, an order of magnitude with a different
operational definition. The model has a growth phase with
socially competent adults, but it is about half as long as Calhoun's phase~B (13 against 30 weeks), although it holds two thirds
of the births, and O10b, which asks for it, passes in 0.28 of the runs when the adult mean is replaced by that of the animals
past the window (0.02 with the adult mean, which the newborns dilute). The phases appear in Calhoun's order, but B
and C are compressed and an artificial plateau precedes D.
Within this transient, crowding damage is the trigger (without it colonies explode), social learning (R2) the
amplifier, the plasticity window (R1) the ratchet and the refuges (R7) make the isolated class persistent. The
ratchet matters because of the refuges: they are pockets of low density in which, without R1, deteriorated animals
recover and breed again (natality resumes and O17p falls to 0.21), whereas without refuges the collapse stays
irreversible in most runs even without R1. R1 is what lets a persistent isolated class coexist with an irreversible
collapse (Section~\ref{sec:model-rules}). Only where
socialisation is subcritical does the learning cascade propagate over several generations and kill a colony without
crowding (Q7), and those colonies first overshoot the crowding threshold, so this is not Calhoun's collapse.

\paragraph{The timing of extinction is extreme-value statistics.}
A post hoc analysis of the runs of the preregistered design suggested that crowd aversion makes the extinction time
more reproducible. It is not a
signature of the mechanism. Without social mortality $T_{\rm ext}\approx t_{\rm last}+\ell_{\max}$, the last birth
plus the longest residual lifetime (Appendix~\ref{app:calib-evs}). The second term, about 110 weeks, fluctuates only
on the scale of the inverse slope of the senescent hazard, as expected for the maximum of many lifetimes
\cite{Gumbel1958,Coles2001}, so the spread of $T_{\rm ext}$ is that of the last birth (correlation 0.92), set by how
abruptly natality stops. Without attraction the spread is 0.04, smaller than in \textsf{C1}, and it falls with the
size of the lattice (0.14 at $L=20$, 0.08 at $L=30$--60; Table~\ref{tab:res-size}).

\paragraph{What ``minimal'' means here.}
``Minimal'' does not mean few parameters: the core model has more than a dozen, and the added rules about ten (R1: $a_w$;
R2: $\lambda$, $w$, $s_{\rm base}$; AV: $\kappa$, $p$; R7: $f_R$, $c_R$, $\pi$, $q$, $b$ and the hard capacity), all
fixed by qualitative calibration against the evaluated patterns. It means that every added rule is local, has a
neutral value that recovers the core model and is motivated by an observation in Calhoun's report; that rules
found unnecessary were removed (long-range jumps, state-dependent attraction, territorial exclusion, social
mortality); and that each rule of the minimal set is necessary by ablation \emph{at the evaluation thresholds used}:
a class threshold of 10\%, isolation at $m\le2$, persistence of four weeks and an ensemble threshold of 0.8.
The necessity of R7 is necessity for O13p, a criterion added together with R7; for O6, O7 and O11 removing R7 changes
nothing, whereas R1 and R2 remain needed (O11 in 0.71 of the applicable runs and 0.00 without each).
The preregistered necessity test Q2 is rejected with the primary criteria ($p_{\rm Holm}=2.2\times10^{-7}$) but not
with the preregistered criteria ($p_{\rm Holm}=0.067$); the rejection comes only from the post hoc choice of $t_{B99}$ in
O5, a generic criterion (Table~\ref{tab:deviations}). Removing AV lowers the primary outcome by 0.20 [0.13, 0.28] in
the factorial. The predicted size of the effect, however, fails: without AV the primary outcome was to fall to at
most 0.5, and it is 0.79 [0.71, 0.85]. AV therefore contributes but is not necessary in the preregistered sense; we
do not use the observed power of the test, which cannot explain a decision \cite{HoenigHeisey2001}.
At these thresholds the minimal set for the primary outcome is \textsf{C1}$'$ (R1, R2 and R7): 0.835 [0.78, 0.88] of
200 runs, borderline because the lower bound is below 0.8, against 0.995 for \textsf{C1} (not preregistered). It falls
to 0.33 at a class threshold of 15\%, and that of
\textsf{C2} from 0.70 to 0.01 at 12.5\% (Fig.~\ref{fig:res-theta}), so minimality is not a property of the rules
alone. Adding AV makes O5 more reliable and the decline senescent (O8 in 0.62 against 0.21), and reduces the free
space at the peak ($f_0=0.16$ against 0.52); partial inheritance is not necessary
($w=0$ with $r=0.15$ gives 1.00). Making the list of required patterns and their thresholds explicit, and testable
run by run, is in our view the main methodological contribution of this work.

\subsection{Limitations and circularity}
\label{sec:limitations}

\paragraph{Circularity between the patterns and the calibration.}
The same patterns and evaluator were used to \emph{choose} the rules and to \emph{report} them; no pattern was held
out. The refuge capacity $c_R=2$ coincides with the isolation threshold of O13 (with $c_R=3$, O17p is zero at every
point of the sensitivity designs); the persistence criterion O13p was added after \textsf{C0}, and R7 was designed to
satisfy it; $m_c$ is both the damage threshold and, in the preregistered evaluator, the boundary of the crowded class (with a fixed
threshold $m>8$ the Calhoun-like fraction of the hypercubes is 2.0\%, against 2.5\% with $m>m_c$). The empty fraction
at the peak was used informally to discard candidates and is a condition of the
primary O4; O12 was used to choose $\alpha=4$ (Section~\ref{sec:circularity}).
The fulfilment of the patterns by \textsf{C1} therefore shows that they are \emph{possible} under local rules, not
that they were predicted. Among the preregistered tests, Q1 and Q4 are consistency checks, the contrast of Q3 is
confounded by the demography, the prediction of Q6 is not supported (with the primary O11 the window has little or no
effect on the boundary), and Q7 establishes an endogenous collapse without crowding, not the form of its boundary. All deviations from the plan are listed in Table~\ref{tab:deviations}.

\paragraph{Robustness.}
Against stochasticity robustness is high: O17p is 0.995 [0.97, 1.00] over 200 runs, independent ensembles replicate
it, eight refuge layouts gave no difference, and exploratory runs on lattices up to $L=84$ at fixed
density (30 runs per size, single-pass refuge admission) give the same peak density and persistent class, with no healthy pocket reseeding the colony. Locally it holds over most of
the phase planes, whose boundaries are estimated from 30 runs per point. Globally it is limited and relative to the
sampled box: over two 15-factor Latin hypercubes the mean primary outcome is 0.040 [0.026, 0.059] and 2.0\%
[0.7, 3.6] of the points are Calhoun-like (7.2\% with the preregistered criteria). Part of this is definitional ($c_R=3$,
large $m_c$) and part is mechanism (a total refuge capacity $f_Rc_R\ge0.3$ and a strong preference $\pi\ge100$); even
with the four restrictions together, chosen after the runs, 33 of 37 points fail. The Sobol' indices are conditional
on the factors fixed after screening. Demographically, the primary outcome depends on the timing of the end of natality, as discussed above, while the
patterns that the rules add do not.

\paragraph{Scales and the comparison with Universe~25.}
With at most one Moore step per week, the Damköhler number $t_{\rm mix}/t_{\rm det}$ is about $1.7\times10^2$ in
\textsf{C1} (about 10 for founder colonies), against $\lesssim4\times10^{-3}$ in Universe~25, whose mice crossed the
enclosure every day (Table~\ref{tab:calib-scales}). The spatial structure of the model is that of animals almost
sessile on the scale of the deterioration, and with more movement sub-steps the coexistence of classes is lost. The
model lacks the separation $t_{\rm mix}\ll T_{\rm gen}\ll t_{\rm det}$ of Universe~25. We map a cell to a nest box, $m_c$ to its comfortable occupancy (``fifteen
mice could comfortably rest in a single nest box'' \cite{Calhoun1973}), $K_\gamma$ to Calhoun's 3840 places and the
refuges to the little-preferred high nests; $c_R=2$ is the isolation threshold, not the capacity of a nest. With this
mapping $f_0$ and $\rho_{\rm pk}$ have the same definitions as Calhoun's 20\% of empty nest boxes and 0.57, but $|G|$
and $m_c$ were not calibrated and $\rho_{\rm pk}$ depends on the seeding density, so they agree only in sign. R7 was
calibrated with synchronous moves; the exact hard capacity leaves the primary outcome unchanged, and with a random
sequential update a soft capacity suffices ($\rho_{\rm pk}=0.43$).

\paragraph{What is not reproduced.}
The founder phase (O1--O3): from four founding pairs a colony spreads as a slow front, and no point of an exploratory
founder plane is Calhoun-like (O17p at most 0.37). The beautiful ones (O13b): the persistent isolated class of
\textsf{C1} is made of refugees of the early boom, most of whom entered the refuges as pups, not of a late-born
cohort; the variant \textsf{C2} produces a later-born class at a demographic cost. Total infant mortality (O15) and the
sex bias of the remnant (O9) are not reproduced, because conception and pup survival decline together and there is no
sex-specific mortality. Details, the refuted coexistence band (Q5) and the comparison with
\cite{GwizdallaDus2026} are in Appendix~\ref{app:discussion}. That model reproduces the founder phase and the timing
of the curve, which ours does not; ours shows that local rules, without population-level switches, produce
coexisting classes of deteriorated animals, one of them persistent, on top of an envelope (a low peak with free
space, an end of natality and the failure of the late-born cohorts) that crowding damage and the demography already
give. Both agree that replicate variability is large and that Calhoun's trajectory is one admissible outcome among
many. We make no claim about human populations.

\paragraph{Future work.}
A founder protocol with realistic mobility and an explicit adaptation phase; patterns held out from the calibration;
gestation, lactation and background mortality, so that births spread up to the peak, which is the condition for a
late-born isolated cohort and for natality to end at high population without a synchronous start; decoupling
conception from pup survival and adding sex-specific risks; and a closer comparison with maternal-effect theories of
population cycles \cite{GinzburgTaneyhill1994,InchaustiGinzburg1998}.
